\chapter{Theoretical Foundations \& Literature Survey}
\label{chap:literature_survey}

\section{The Theoretical Evolution of Option Pricing}
\label{sec:option_pricing_evolution}

The modern quantitative modeling of financial derivatives originates from the foundational work of Fischer Black, Myron Scholes \cite{black1973pricing}, and Robert Merton \cite{merton1973theory}. The core insight of the Black-Scholes-Merton (BSM) framework was the construction of a continuously rebalanced, risk-free hedging portfolio combining the underlying asset and a derivative contract. By eliminating market risk through continuous dynamic hedging, Black, Scholes, and Merton derived a closed-form analytical solution for European option prices.

For a European call option $C$ and put option $P$ with strike price $K$ and time-to-expiry $\tau = T - t$, the canonical BSM formulas are:
\begin{equation}
    C = S_t N(d_1) - K e^{-r\tau} N(d_2), \quad P = K e^{-r\tau} N(-d_2) - S_t N(-d_1)
    \label{eq:bsm_call_lit}
\end{equation}
where $S_t$ is the spot price of the underlying asset, $r$ is the continuously compounded risk-free interest rate, $N(\cdot)$ is the cumulative standard normal distribution function, and:
\begin{equation}
    d_1 = \frac{\ln(S_t / K) + \left(r + \frac{1}{2}\sigma^2\right)\tau}{\sigma \sqrt{\tau}}, \quad d_2 = d_1 - \sigma \sqrt{\tau}
\end{equation}

The model relies on several simplifying market assumptions:
\begin{itemize}\setlength{\itemsep}{0pt}\setlength{\parskip}{0pt}
    \item \textbf{Geometric Brownian Motion:} Asset prices follow continuous diffusion with constant drift and constant return volatility $\sigma$.
    \item \textbf{Frictionless Electronic Markets:} Zero transaction costs, zero bid-ask spread friction, and unrestricted short-selling.
    \item \textbf{Continuous Trading:} Portfolios can be continuously rebalanced at infinitesimal intervals without price slippage or market impact.
    \item \textbf{Stationary Parameters:} The risk-free rate $r$ and volatility $\sigma$ remain known and constant throughout the life of the contract.
\end{itemize}
While the BSM formula remains the standard quoting convention of global derivatives exchanges, its theoretical assumption of constant volatility is systematically violated in live financial markets.

\section{The Implied Volatility Surface: Smile, Skew, and Term Structure}
\label{sec:iv_surface_dynamics}

In live options markets, volatility $\sigma$ is an unobservable parameter. Traders invert the BSM formula numerically using observed transaction prices $C_{\text{market}}$ to determine the \textbf{Implied Volatility ($\sigma_{\text{implied}}$)}:
\begin{equation}
    C_{\text{market}} - C_{\text{BSM}}(S_t, K, \tau, r, \sigma_{\text{implied}}) = 0
\end{equation}

If the BSM model reflected reality perfectly, inverting the formula across different strikes $K$ and expiration dates $T$ would yield a constant volatility plane. However, empirical option markets exhibit distinct structural variations across strike and maturity, giving rise to the 3D \textbf{Implied Volatility Surface}:
\begin{itemize}\setlength{\itemsep}{0pt}\setlength{\parskip}{0pt}
    \item \textbf{Pronounced Negative Skew:} Out-of-the-money (OTM) put options (low strikes, $K < S$) trade at significantly higher implied volatilities than out-of-the-money call options, reflecting downside portfolio insurance demand ("crashophobia") and corporate leverage.
    \item \textbf{Smile Curvature (Convexity):} OTM and deep ITM wings exhibit elevated implied volatility relative to at-the-money (ATM) contracts, pricing fat-tailed return distributions.
    \item \textbf{Term Structure of Volatility:} Implied volatility varies systematically with time-to-expiry, sloping upward in calm regimes (contango) and spiking during market stress (backwardation).
\end{itemize}

To resolve the volatility smile, continuous-time finance introduced \textbf{Local Volatility} models (Dupire, 1994; Derman and Kani \cite{derman1994riding}) and \textbf{Stochastic Volatility} models (most notably Heston \cite{heston1993closed}). While stochastic volatility allows volatility itself to follow a mean-reverting random process correlated with the underlying spot price, these continuous-time models require intensive computational recalibration and cannot directly incorporate discrete limit order book microstructure features.

\section{Econometric Approaches to Volatility Modeling}
\label{sec:econometric_approaches}

Before the adoption of machine learning, financial econometricians modeled time-series volatility through discrete-time autoregressive processes, pioneered by Robert Engle's \textbf{ARCH} model \cite{engle1982autoregressive} and generalized by Tim Bollerslev into \textbf{GARCH} \cite{bollerslev1986generalized}. In high-frequency finance, Andersen et al. \cite{andersen2003modeling} showed that summing intraday squared returns yields non-parametric measures of \textbf{Realized Volatility (RV)}.

While ARCH, GARCH, and Realized Volatility capture volatility clustering and persistence in historical stock returns, they model the physical probability measure ($\mathbb{P}$) of the underlying asset rather than the forward-looking risk-neutral measure ($\mathbb{Q}$) priced into option contracts. Furthermore, classical time-series models assume parametric linear relationships that struggle to adapt to the ultra-short, non-linear 5-minute sampling intervals of modern electronic derivatives desks.

\section{Machine Learning \& Non-Parametric Ensembles in Derivatives}
\label{sec:ml_in_derivatives}

Over the past decade, institutional trading firms have increasingly adopted non-parametric machine learning architectures for derivatives modeling \cite{lopez2018advances}. Machine learning provides three distinct advantages over classical econometric models:
\begin{enumerate}\setlength{\itemsep}{0pt}\setlength{\parskip}{0pt}
    \item \textbf{Model-Free Flexibility:} Algorithms learn complex volatility dynamics directly from market data without imposing rigid diffusion assumptions.
    \item \textbf{Multi-Dimensional Feature Ingestion:} Ensembles can simultaneously process cross-sectional strike geometry, multi-horizon price momentum, and order-book liquidity metrics.
    \item \textbf{Non-Linear Interaction Capture:} Tree-based models automatically capture non-linear interactions, such as how the impact of order-book spreads on volatility changes depending on whether a contract is at-the-money or deep out-of-the-money.
\end{enumerate}

\subsection{Random Forests}
Introduced by Leo Breiman in 2001 \cite{breiman2001random}, \textbf{Random Forests} construct an ensemble of decorrelated decision trees using bootstrap aggregation (bagging) and random feature subspace selection. Each individual decision tree recursively partitions the feature space into homogeneous regions to minimize prediction variance. In high-frequency options forecasting, Random Forests excel because averaging across hundreds of decorrelated trees dampens idiosyncratic quote noise and avoids overfitting to temporary bid-ask bounces.

\subsection{Gradient Boosted Decision Trees (XGBoost)}
In contrast to bagging, **Gradient Boosting** builds trees sequentially in an additive stage-wise manner. Proposed by Chen and Guestrin in 2016 \cite{chen2016xgboost}, \textbf{XGBoost} (eXtreme Gradient Boosting) optimizes a regularized objective function combining prediction loss with tree-complexity penalties. By utilizing first- and second-order loss gradients, XGBoost fits residuals with high computational efficiency, providing a powerful non-linear alternative to Random Forests.

\section{Market Microstructure, Limit Order Books, and Intraday Liquidity}
\label{sec:market_microstructure}

Prices on modern electronic exchanges are formed through continuous double-auction mechanisms in \textbf{Limit Order Books (LOB)} \cite{hasbrouck2007empirical, bouchaud2003theory}. At any discrete timestamp $t$, the state of an option contract is defined by standing limit buy orders (bids) and limit sell orders (asks). The \textbf{Bid-Ask Spread} ($\text{Spread}_t = \text{Ask}_t - \text{Bid}_t$) represents the instantaneous cost of liquidity immediacy.

In market microstructure theory, the bid-ask spread reflects three primary operational frictions:
\begin{itemize}\setlength{\itemsep}{0pt}\setlength{\parskip}{0pt}
    \item \textbf{Order Processing Costs:} Direct exchange routing, clearing, and regulatory fees incurred by market makers.
    \item \textbf{Inventory Holding Risk:} The risk of holding unhedged options positions during sudden market reversals.
    \item \textbf{Adverse Selection (Informed Trading):} As formulated by Glosten and Milgrom \cite{glosten1985bid}, market makers widen spreads when facing informed counterparties possessing superior information about impending underlying moves or volatility shifts.
\end{itemize}
In intraday options trading, sudden spread widening frequently precedes sharp adjustments in implied volatility. Incorporating order book spreads directly into feature engineering provides machine learning models with early warning signals of liquidity withdrawal. Furthermore, near weekly settlement, market makers face **Pin Risk**—the uncertainty over whether an at-the-money option will expire in- or out-of-the-money—leading to extreme order book volatility.

\section{Identification of the Research Gap}
\label{sec:research_gap}

A review of the derivatives literature highlights four critical gaps that this study directly addresses:
\begin{enumerate}\setlength{\itemsep}{0pt}\setlength{\parskip}{0pt}
    \item \textbf{Scarcity of Ultra-High-Frequency (5-Minute) IV Studies:} Most published volatility literature focuses on daily or monthly horizons; empirical models forecasting 5-minute implied volatility changes across cross-sectional option chains in emerging markets remain exceptionally rare.
    \item \textbf{Unrealistic Benchmark Evaluations:} Academic papers frequently evaluate machine learning models using randomized cross-validation that leaks future temporal data, and rarely benchmark against the formidable **Persistence Random Walk** ($\widehat{\Delta \text{IV}}_{t+5} = 0$).
    \item \textbf{Lack of Operational Valuation Translation:} Machine learning studies often terminate at statistical error metrics (MAE/RMSE) without converting predictions into economic signals (such as relative valuation screeners or scenario stress grids).
    \item \textbf{Neglect of the Expiry-Day Boundary Discontinuity:} Literature regularly assumes stationary dynamics throughout the contract lifecycle, neglecting the severe structural collapse occurring in the hyper-gamma, vanishing-vega final hours of weekly expiry.
\end{enumerate}

This study bridges these gaps by developing a complete, application-oriented quantitative system combining automated Upstox API data ingestion, strict chronological benchmarking, conviction-gated trading rules, relative valuation screener mechanics, and a forensic analysis of the expiry-day regime collapse.
